\documentclass[11pt]{article}
\usepackage[T1]{fontenc}
\usepackage[utf8]{inputenc}
\usepackage{lmodern}
\usepackage[margin=1in]{geometry}
\usepackage{amsmath,amssymb}
\usepackage{booktabs,tabularx}
\usepackage{enumitem}
\usepackage{hyperref}

\IfFileExists{paper/release_info.tex}{
  \newcommand{\OPHPaperReleaseID}{r2040}
\newcommand{\OPHPaperReleaseDate}{September 9, 2026}
\newcommand{\OPHPaperReleaseBanner}{%
\begin{center}
\small\textbf{Paper release:} \texttt{\OPHPaperReleaseID}\quad\textbf{Released:} \OPHPaperReleaseDate
\end{center}
}
\newcommand{\PaperReleaseID}{\OPHPaperReleaseID}
\newcommand{\PaperReleaseDate}{\OPHPaperReleaseDate}
\newcommand{\PaperReleaseBanner}{%
\begin{center}
\small\textbf{Paper release:} \texttt{\PaperReleaseID}\quad\textbf{Released:} \PaperReleaseDate
\end{center}
}

}{
  \newcommand{\OPHPaperReleaseID}{r2040}
\newcommand{\OPHPaperReleaseDate}{September 9, 2026}
\newcommand{\OPHPaperReleaseBanner}{%
\begin{center}
\small\textbf{Paper release:} \texttt{\OPHPaperReleaseID}\quad\textbf{Released:} \OPHPaperReleaseDate
\end{center}
}
\newcommand{\PaperReleaseID}{\OPHPaperReleaseID}
\newcommand{\PaperReleaseDate}{\OPHPaperReleaseDate}
\newcommand{\PaperReleaseBanner}{%
\begin{center}
\small\textbf{Paper release:} \texttt{\PaperReleaseID}\quad\textbf{Released:} \PaperReleaseDate
\end{center}
}

}

\hypersetup{colorlinks=true,linkcolor=blue,citecolor=blue,urlcolor=blue}
\setlength{\parindent}{0pt}
\setlength{\parskip}{0.55em}
\setlength{\emergencystretch}{3em}

\title{Observer-Patch Holography Cosmology Data and Likelihood Contracts\\
\large Provenance, Comparability, and Claim Boundaries}
\author{B.\ M\"uller}
\date{\OPHPaperReleaseDate}

\begin{document}
\maketitle
\OPHPaperReleaseBanner

\begin{center}
\small\textbf{Author affiliation:} Bernhard Mueller, Pragma Research Inc.
\end{center}

\begin{abstract}
We specify data and likelihood contracts for conditional Observer-Patch
Holography (OPH) cosmology. Each comparison binds a finite source, compatible
transfer equations, declared datasets and covariance, nuisance parameters,
and a common observable projection. The contracts distinguish source-derived
quantities, imported inputs and fitted parameters, and account for sample
overlap and shared calibration. They give a reproducible interface from
screen and dark-sector models to cosmological data.
\end{abstract}

\section{Comparison Classes}

Every reported number belongs to one of four classes:
\begin{enumerate}[leftmargin=2em]
\item \textbf{source-derived}: all physical inputs follow from the declared OPH
construction without using the comparison data;
\item \textbf{imported}: established background or microphysical inputs are
declared explicitly;
\item \textbf{fit}: one or more parameters are inferred from the same data;
\item \textbf{diagnostic}: the calculation tests shape, scale, code behavior, or
an interface without supporting a physical OPH claim.
\end{enumerate}

The cosmology comparisons in scope are diagnostic or imported. No source-derived
joint cosmology likelihood exists.

The exact finite source result computes source height from authenticated
direct-parent edges as the exact longest parent-chain length, attained at
every event, and proves strict increase along generated ancestry. Independently, a real axis and the rank-three source quotient define
a one-plus-three ambient Lorentzian precursor. An explicit positive
\texttt{timeScale} on source height enters only through event placement. This
construction also gives every finite log an auxiliary enumerated injective
forward-causal placement, without order reflection or physical coordinate
selection. It does not determine intrinsic poset dimension or supply a
physical clock or continuum, background solution, conserved cosmological
stress, initial conditions, transfer functions, or an observable likelihood.
Those objects remain subject to the model-information and no-data-use rules
below.

The \href{https://github.com/FloatingPragma/observer-patch-holography/blob/main/paper/recovering_observer_spacetime_and_einstein_dynamics_from_overlap_consistency.tex}{spacetime construction} gives a flat causal and normalized count-volume limit for
conservative record populations with a specified neighbour-read law and
model clock. In a likelihood comparison, the selection and physical
calibration of that law count as model information, alongside the
cosmological background and matter assumptions.

\section{Required Model Information}

A physical comparison must state:
\begin{enumerate}[leftmargin=2em]
\item the covariant action and source map;
\item the background solution and species content;
\item the gauge-invariant initial conditions;
\item the perturbation, collision, and exchange equations;
\item the observable projection and numerical tolerances;
\item every external input and fitted parameter.
\end{enumerate}

For the modular-charge dark-sector proposal, conserved charge in a supplied
comoving volume gives the homogeneous dilution law \(\rho_A\propto a^{-3}\).
With the separate no-exchange continuity equation it gives a pressureless
background. The abundance, relativistic stress,
initial perturbations, lensing map, clusters, and nonlinear structure equations
are not supplied here.

\section{Required Data Information}

A comparison must identify the data release, selection, masks, covariance,
calibration model, nuisance parameters, priors, and combination rule. Dataset
overlap and shared calibration must be included. A diagonal sum is not a joint
likelihood when the covariance is non-diagonal.

The model and data layers are separate. Measured values that select
a source formula, normalization, branch, or applicability domain are inputs to
that comparison rather than OPH outputs.

\section{Observable Contracts}

\begin{center}
\begin{tabularx}{\textwidth}{@{}l>{\raggedright\arraybackslash}X>{\raggedright\arraybackslash}X@{}}
\toprule
Observable family & Required theory output & Required comparison object \\
\midrule
Cosmic microwave background (CMB) & $C_\ell^{TT}$, $C_\ell^{TE}$, $C_\ell^{EE}$, lensing, foreground model &
map or band-power likelihood with covariance and nuisance model \\
Baryon acoustic oscillations (BAO) and supernovae & $H(z)$, $D_A(z)$, luminosity distance, sound horizon &
survey likelihood with calibration covariance \\
Weak lensing & metric potentials, nonlinear matter power, intrinsic alignments &
tomographic likelihood with masks and covariance \\
Growth and redshift-space distortions (RSD) & $P_m(k,z)$ and $f\sigma_8(z)$ &
windowed survey likelihood with cross-bin covariance \\
Clusters & mass function, selection, observable--mass map, lensing calibration &
count and calibration likelihood \\
Galaxies & disk dynamics, gas and stellar nuisances, external environment &
galaxy-level likelihood with selection and covariance \\
\bottomrule
\end{tabularx}
\end{center}

\section{Diagnostic Results}

The reported CMB overlays, compressed background comparisons, and galaxy
scaling checks are diagnostic. Their formulas or inputs are selected with knowledge of
the comparison data, omit required covariance or nuisance structure, or lack a
physical source map. They test numerical paths and conditional formulas and
do not establish a cosmological prediction.

\section{No-Data-Use Boundary}

A source-derived calculation must expose its dependency graph. The source side
may use mathematical constants, declared OPH finite artifacts, and explicitly
imported physical inputs. It may not use the observed quantity being claimed as
an output, a fit to that quantity, or a proxy calibrated from the same dataset.

This provenance rule is enforced through the declared dependency graph and
data-use classification.

\section{Theorem scope}

The screen-spectrum and radial-lift mathematics are conditional theorems.
One-shell inversion is nonunique; physical source dilation and radial
cross-covariance tomography are two established sufficient uniqueness routes.
Finite source instantiation, repair-charge cosmological source, Boltzmann
bridge, and joint likelihood are not supplied. All cosmological
data products are diagnostic. No cosmological result carries an OPH
falsification verdict.

\section*{AI Assistance Disclosure}
This research project used research-grade commercial models, including
Anthropic's Fable and OpenAI's GPT-5.6-Sol, for research support, software
development, editing, and synthesis. The authors are responsible for the
paper's claims, methods, and final text.

\end{document}
